\section{Filtrations}
\label{section-filtrations}

\noindent
A nice reference for this material is \cite[Section 1]{HodgeII}.
(Note that our conventions regarding abelian categories are different.)

\begin{definition}
\label{definition-filtered}
Let $\mathcal{A}$ be an abelian category.
\begin{enumerate}
\item A {\it decreasing filtration} $F$ on an object $A$
is a family $(F^nA)_{n \in \mathbf{Z}}$ of subobjects of $A$ such that
$$
A \supset \ldots \supset F^nA \supset F^{n + 1}A \supset \ldots \supset 0
$$
\item A {\it filtered object of $\mathcal{A}$} is
pair $(A, F)$ consisting of an object $A$ of $\mathcal{A}$
and a decreasing filtration $F$ on $A$.
\item A {\it morphism $(A, F) \to (B, F)$ of filtered objects}
is given by a morphism $\varphi : A \to B$ of $\mathcal{A}$
such that $\varphi(F^iA) \subset F^iB$ for all $i \in \mathbf{Z}$.
\item The category of filtered objects is denoted $\text{Fil}(\mathcal{A})$.
\item Given a filtered object $(A, F)$ and a subobject $X \subset A$ the
{\it induced filtration} on $X$ is the filtration with $F^nX = X \cap F^nA$.
\item Given a filtered object $(A, F)$ and a surjection
$\pi : A \to Y$ the {\it quotient filtration} is the filtration with
$F^nY = \pi(F^nA)$.
\item A filtration $F$ on an object $A$ is said to be {\it finite}
if there exist $n, m$ such that $F^nA = A$ and $F^mA = 0$.
\item Given a filtered object $(A, F)$ we say $\bigcap F^iA$ exists
if there exists a biggest subobject of $A$ contained in all $F^iA$.
We say $\bigcup F^iA$ exists if there exists a smallest subobject
of $A$ containing all $F^iA$.
\item The filtration on a filtered object $(A, F)$ is said to be
{\it separated} if $\bigcap F^iA = 0$ and
{\it exhaustive} if $\bigcup F^iA = A$.
\end{enumerate}
\end{definition}

\noindent
By abuse of notation we say that a morphism $f : (A, F) \to (B, F)$
of filtered objects is {\it injective} if $f : A \to B$ is injective
in the abelian category $\mathcal{A}$. Similarly we say $f$ is
{\it surjective} if $f : A \to B$ is surjective in the category
$\mathcal{A}$. Being injective (resp.\ surjective)
is equivalent to being a monomorphism (resp.\ epimorphism)
in $\text{Fil}(\mathcal{A})$. By
Lemma \ref{lemma-filtered}
this is also equivalent to having zero kernel (resp.\ cokernel).

\begin{lemma}
\label{lemma-filtered}
Let $\mathcal{A}$ be an abelian category.
The category of filtered objects $\text{Fil}(\mathcal{A})$
has the following properties:
\begin{enumerate}
\item It is an additive category.
\item It has a zero object.
\item It has kernels and cokernels, images and coimages.
\item In general it is not an abelian category.
\end{enumerate}
\end{lemma}

\begin{proof}
It is clear that $\text{Fil}(\mathcal{A})$ is additive with direct
sum given by $(A, F) \oplus (B, F) = (A \oplus B, F)$ where
$F^p(A \oplus B) = F^pA \oplus F^pB$.
The kernel of a morphism $f : (A, F) \to (B, F)$ of filtered
objects is the injection $\Ker(f) \subset A$ where $\Ker(f)$
is endowed with the induced filtration.
The cokernel of a morphism $f : A \to B$ of filtered
objects is the surjection $B \to \Coker(f)$ where $\Coker(f)$
is endowed with the quotient filtration. Since all kernels and cokernels
exist, so do all coimages and images. See
Example \ref{example-not-abelian}
for the last statement.
\end{proof}

\begin{definition}
\label{definition-strict}
Let $\mathcal{A}$ be an abelian category.
A morphism $f : A \to B$ of filtered objects of $\mathcal{A}$ is
said to be {\it strict} if $f(F^iA) = f(A) \cap F^iB$ for
all $i \in \mathbf{Z}$.
\end{definition}

\noindent
This also equivalent to requiring that $f^{-1}(F^iB) = F^iA + \Ker(f)$
for all $i \in \mathbf{Z}$. We characterize strict morphisms
as follows.

\begin{lemma}
\label{lemma-characterize-strict-general}
Let $\mathcal{A}$ be an abelian category.
Let $f : A \to B$ be a morphism of filtered objects of $\mathcal{A}$.
The following are equivalent
\begin{enumerate}
\item $f$ is strict,
\item the morphism $\Coim(f) \to \Im(f)$ of
Lemma \ref{lemma-coim-im-map}
is an isomorphism.
\end{enumerate}
\end{lemma}

\begin{proof}
Note that $\Coim(f) \to \Im(f)$ is an isomorphism of
objects of $\mathcal{A}$, and that part (2) signifies that it is
an isomorphism of filtered objects.
By the description of kernels and cokernels in the proof of
Lemma \ref{lemma-filtered}
we see that the filtration on $\Coim(f)$ is the
quotient filtration coming from $A \to \Coim(f)$.
Similarly, the filtration on $\Im(f)$ is the induced
filtration coming from the injection $\Im(f) \to B$.
The definition of strict is exactly that the quotient filtration
is the induced filtration.
\end{proof}

\begin{lemma}
\label{lemma-add-summand-strict-monomorphism}
Let $\mathcal{A}$ be an abelian category.
Let $f : A \to B$ be a strict monomorphism of filtered objects.
Let $g : A \to C$ be a morphism of filtered objects.
Then $f \oplus g : A \to B \oplus C$ is a strict monomorphism.
\end{lemma}

\begin{proof}
Clear from the definitions.
\end{proof}

\begin{lemma}
\label{lemma-add-summand-strict-epimorphism}
Let $\mathcal{A}$ be an abelian category.
Let $f : B \to A$ be a strict epimorphism of filtered objects.
Let $g : C \to A$ be a morphism of filtered objects.
Then $f \oplus g : B \oplus C \to A$ is a strict epimorphism.
\end{lemma}

\begin{proof}
Clear from the definitions.
\end{proof}

\begin{lemma}
\label{lemma-induced-and-quotient-strict}
Let $\mathcal{A}$ be an abelian category.
Let $(A, F)$, $(B, F)$ be filtered objects.
Let $u : A \to B$ be a morphism of filtered objects.
If $u$ is injective then $u$ is strict if and only if the filtration
on $A$ is the induced filtration.
If $u$ is surjective then $u$ is strict if and only if the filtration
on $B$ is the quotient filtration.
\end{lemma}

\begin{proof}
This is immediate from the definition.
\end{proof}

\begin{lemma}
\label{lemma-composition-strict}
Let $\mathcal{A}$ be an abelian category. Let $f : A \to B$, $g : B \to C$
be strict morphisms of filtered objects.
\begin{enumerate}
\item In general the composition $g \circ f$ is not strict.
\item If $g$ is injective, then $g \circ f$ is strict.
\item If $f$ is surjective, then $g \circ f$ is strict.
\end{enumerate}
\end{lemma}

\begin{proof}
Let $B$ a vector space over a field $k$ with basis $e_1, e_2$, with the
filtration $F^nB = B$ for $n < 0$, with $F^0B = ke_1$, and $F^nB = 0$ for
$n > 0$. Now take $A = k(e_1 + e_2)$ and $C = B/ke_2$ with filtrations
induced by $B$, i.e., such that $A \to B$ and $B \to C$ are strict
(Lemma \ref{lemma-induced-and-quotient-strict}).
Then $F^n(A) = A$ for $n < 0$ and $F^n(A) = 0$ for $n \geq 0$. 
Also $F^n(C) = C$ for $n \leq 0$ and $F^n(C) = 0$ for $n > 0$.
So the (nonzero) composition $A \to C$ is not strict.

\medskip\noindent
Assume $g$ is injective. Then
\begin{align*}
g(f(F^pA)) & = g(f(A) \cap F^pB) \\
& = g(f(A)) \cap g(F^p(B)) \\
& = (g \circ f)(A) \cap (g(B) \cap F^pC) \\
& = (g \circ f)(A) \cap F^pC.
\end{align*}
The first equality as $f$ is strict, the second because $g$ is injective,
the third because $g$ is strict, and the fourth because
$(g \circ f)(A) \subset g(B)$.

\medskip\noindent
Assume $f$ is surjective. Then
\begin{align*}
(g \circ f)^{-1}(F^iC) & = f^{-1}(F^iB + \Ker(g)) \\
& = f^{-1}(F^iB) + f^{-1}(\Ker(g)) \\
& = F^iA + \Ker(f) + \Ker(g \circ f) \\
& = F^iA + \Ker(g \circ f)
\end{align*}
The first equality because $g$ is strict, the second because $f$ is
surjective, the third because $f$ is strict, and the last because
$\Ker(f) \subset \Ker(g \circ f)$.
\end{proof}

\noindent
The following lemma says that subobjects of a filtered object have a well
defined filtration independent of a choice of writing the object as a
cokernel.

\begin{lemma}
\label{lemma-filtration-subobject}
Let $\mathcal{A}$ be an abelian category.
Let $(A, F)$ be a filtered object of $\mathcal{A}$.
Let $X \subset Y \subset A$ be subobjects of $A$.
On the object
$$
Y/X = \Ker(A/X \to A/Y)
$$
the quotient filtration coming from the induced filtration on $Y$ and the
induced filtration coming from the quotient filtration on $A/X$ agree.
Any of the morphisms $X \to Y$, $X \to A$, $Y \to A$, $Y \to A/X$,
$Y \to Y/X$, $Y/X \to A/X$ are strict (with induced/quotient filtrations).
\end{lemma}

\begin{proof}
The quotient filtration $Y/X$ is given by
$F^p(Y/X) = F^pY/(X \cap F^pY) = F^pY/F^pX$
because $F^pY = Y \cap F^pA$ and $F^pX = X \cap F^pA$.
The induced filtration from the injection $Y/X \to A/X$ is given by
\begin{align*}
F^p(Y/X) & = Y/X \cap F^p(A/X) \\
& = Y/X \cap (F^pA + X)/X \\
& = (Y \cap F^pA)/(X \cap F^pA) \\
& = F^pY/F^pX.
\end{align*}
Hence the first statement of the lemma.
The proof of the other cases is similar.
\end{proof}

\begin{lemma}
\label{lemma-pushout-filtered}
Let $\mathcal{A}$ be an abelian category.
Let $A, B, C \in \text{Fil}(\mathcal{A})$.
Let $f : A \to B$ and $g : A \to C$ be morphisms.
Then there exists a pushout
$$
\xymatrix{
A \ar[r]_f \ar[d]_g & B \ar[d]^{g'} \\
C \ar[r]^{f'} & C \amalg_A B
}
$$
in $\text{Fil}(\mathcal{A})$. If $f$ is strict, so is $f'$.
\end{lemma}

\begin{proof}
Set $C \amalg_A B$ equal to $\Coker((g, -f) : A \to C \oplus B)$
in $\text{Fil}(\mathcal{A})$. This cokernel exists, by
Lemma \ref{lemma-filtered}.
It is a pushout, see
Example \ref{example-fibre-product-pushouts}.
Note that $F^p(C \amalg_A B)$ is the image of $F^pC \oplus F^pB$.
Hence
$$
(f')^{-1}(F^p(C \amalg_A B)) = g(f^{-1}(F^pB)) + F^pC
$$
Whence the last statement.
\end{proof}

\begin{lemma}
\label{lemma-fibre-product-filtered}
Let $\mathcal{A}$ be an abelian category.
Let $A, B, C \in \text{Fil}(\mathcal{A})$.
Let $f : B \to A$ and $g : C \to A$ be morphisms.
Then there exists a fibre product
$$
\xymatrix{
B \times_A C \ar[r]_{g'} \ar[d]_{f'} & B \ar[d]^f \\
C \ar[r]^g & A
}
$$
in $\text{Fil}(\mathcal{A})$. If $f$ is strict, so is $f'$.
\end{lemma}

\begin{proof}
This lemma is dual to
Lemma \ref{lemma-pushout-filtered}.
\end{proof}

\noindent
Let $\mathcal{A}$ be an abelian category. Let $(A, F)$ be a filtered
object of $\mathcal{A}$. We denote $\text{gr}^p_F(A) = \text{gr}^p(A)$
the object $F^pA/F^{p + 1}A$ of $\mathcal{A}$. This defines an
additive functor
$$
\text{gr}^p :
\text{Fil}(\mathcal{A})
\longrightarrow
\mathcal{A}, \quad
(A, F)
\longmapsto
\text{gr}^p(A).
$$
Recall that we have defined the category $\text{Gr}(\mathcal{A})$
of graded objects of $\mathcal{A}$ in Section \ref{section-graded}.
For $(A, F)$ in $\text{Fil}(\mathcal{A})$ we may set
$$
\text{gr}(A) = \text{the graded object of }\mathcal{A}\text{ whose }
p\text{th graded piece is }\text{gr}^p(A)
$$
and if $\mathcal{A}$ has countable direct sums, then we simply have
$$
\text{gr}(A) = \bigoplus \text{gr}^p(A)
$$
This defines an additive functor
$$
\text{gr} :
\text{Fil}(\mathcal{A})
\longrightarrow
\text{Gr}(\mathcal{A}), \quad
(A, F) \longmapsto \text{gr}(A).
$$

\begin{lemma}
\label{lemma-ses-gr}
Let $\mathcal{A}$ be an abelian category.
\begin{enumerate}
\item Let $A$ be a filtered object and $X \subset A$. Then for each $p$
the sequence
$$
0 \to \text{gr}^p(X) \to \text{gr}^p(A) \to \text{gr}^p(A/X) \to 0
$$
is exact (with induced filtration on $X$ and quotient filtration on $A/X$).
\item Let $f : A \to B$ be a morphism of filtered objects of $\mathcal{A}$.
Then for each $p$ the sequences
$$
0 \to \text{gr}^p(\Ker(f)) \to \text{gr}^p(A) \to
\text{gr}^p(\Coim(f)) \to 0
$$
and
$$
0 \to \text{gr}^p(\Im(f)) \to \text{gr}^p(B) \to
\text{gr}^p(\Coker(f)) \to 0
$$
are exact.
\end{enumerate}
\end{lemma}

\begin{proof}
We have $F^{p + 1}X = X \cap F^{p + 1}A$, hence map
$\text{gr}^p(X) \to \text{gr}^p(A)$ is injective. Dually the map
$\text{gr}^p(A) \to \text{gr}^p(A/X)$ is surjective.
The kernel of $F^pA/F^{p + 1}A \to A/X + F^{p + 1}A$
is clearly $F^{p + 1}A + X \cap F^pA/F^{p + 1}A = F^pX/F^{p + 1}X$
hence exactness in the middle.
The two short exact sequence of (2) are special cases of the
short exact sequence of (1).
\end{proof}

\begin{lemma}
\label{lemma-characterize-strict}
Let $\mathcal{A}$ be an abelian category.
Let $f : A \to B$ be a morphism of finite
filtered objects of $\mathcal{A}$. The following are equivalent
\begin{enumerate}
\item $f$ is strict,
\item the morphism $\Coim(f) \to \Im(f)$ is an isomorphism,
\item $\text{gr}(\Coim(f)) \to \text{gr}(\Im(f))$ is an
isomorphism,
\item the sequence
$\text{gr}(\Ker(f)) \to \text{gr}(A) \to \text{gr}(B)$
is exact,
\item the sequence $\text{gr}(A) \to \text{gr}(B) \to
\text{gr}(\Coker(f))$ is exact, and
\item the sequence
$$
0 \to
\text{gr}(\Ker(f)) \to
\text{gr}(A) \to
\text{gr}(B) \to
\text{gr}(\Coker(f)) \to 0
$$
is exact.
\end{enumerate}
\end{lemma}

\begin{proof}
The equivalence of (1) and (2) is
Lemma \ref{lemma-characterize-strict-general}.
By
Lemma \ref{lemma-ses-gr}
we see that (4), (5), (6) imply (3) and that (3) implies (4), (5), (6).
Hence it suffices to show that (3) implies (2).
Thus we have to show that if $f : A \to B$ is an injective and surjective
map of finite filtered objects which induces and isomorphism
$\text{gr}(A) \to \text{gr}(B)$, then $f$ induces an isomorphism of
filtered objects. In other words, we have to show that
$f(F^pA) = F^pB$ for all $p$.
As the filtrations are finite we may prove this by descending induction
on $p$. Suppose that $f(F^{p + 1}A) = F^{p + 1}B$.
Then commutative diagram
$$
\xymatrix{
0 \ar[r] &
F^{p + 1}A \ar[r] \ar[d]^f &
F^pA \ar[r] \ar[d]^f &
\text{gr}^p(A) \ar[r] \ar[d]^{\text{gr}^p(f)} &
0 \\
0 \ar[r] &
F^{p + 1}B \ar[r] &
F^pB \ar[r] &
\text{gr}^p(B) \ar[r] &
0
}
$$
and the five lemma imply that $f(F^pA) = F^pB$.
\end{proof}

\begin{lemma}
\label{lemma-filtered-complex}
Let $\mathcal{A}$ be an abelian category. Let $A \to B \to C$ be a complex
of filtered objects of $\mathcal{A}$. Assume $\alpha : A \to B$ and
$\beta : B \to C$ are strict morphisms of filtered objects. Then
$\text{gr}(\Ker(\beta)/\Im(\alpha)) =
\Ker(\text{gr}(\beta))/\Im(\text{gr}(\alpha)))$.
\end{lemma}

\begin{proof}
This follows formally from
Lemma \ref{lemma-ses-gr}
and the fact that
$\Coim(\alpha) \cong \Im(\alpha)$ and
$\Coim(\beta) \cong \Im(\beta)$ by
Lemma \ref{lemma-characterize-strict-general}.
\end{proof}

\begin{lemma}
\label{lemma-filtered-acyclic}
Let $\mathcal{A}$ be an abelian category.
Let $A \to B \to C$ be a complex of filtered objects of $\mathcal{A}$.
Assume $A, B, C$ have finite filtrations and that
$\text{gr}(A) \to \text{gr}(B) \to \text{gr}(C)$ is exact.
Then
\begin{enumerate}
\item for each $p \in \mathbf{Z}$ the sequence
$\text{gr}^p(A) \to \text{gr}^p(B) \to \text{gr}^p(C)$ is exact,
\item for each $p \in \mathbf{Z}$ the sequence
$F^p(A) \to F^p(B) \to F^p(C)$ is exact,
\item for each $p \in \mathbf{Z}$ the sequence
$A/F^p(A) \to B/F^p(B) \to C/F^p(C)$ is exact,
\item the maps $A \to B$ and $B \to C$ are strict, and
\item $A \to B \to C$ is exact (as a sequence in $\mathcal{A}$).
\end{enumerate}
\end{lemma}

\begin{proof}
Part (1) is immediate from the definitions.
We will prove (3) by induction on the length of the filtrations.
If each of $A$, $B$, $C$ has only one
nonzero graded part, then (3) holds as $\text{gr}(A) = A$, etc.
Let $n$ be the largest integer such that at least one of
$F^nA, F^nB, F^nC$ is nonzero. Set $A' = A/F^nA$, $B' = B/F^nB$,
$C' = C/F^nC$ with induced filtrations. Note that
$\text{gr}(A) = F^nA \oplus \text{gr}(A')$
and similarly for $B$ and $C$. The induction hypothesis
applies to $A' \to B' \to C'$, which implies that
$A/F^p(A) \to B/F^p(B) \to C/F^p(C)$ is exact for $p \geq n$.
To conclude the same for $p = n + 1$, i.e., to prove that $A \to B \to C$
is exact we use the commutative diagram
$$
\xymatrix{
0 \ar[r] & F^nA \ar[r] \ar[d] & A \ar[r] \ar[d] & A' \ar[r] \ar[d] & 0 \\
0 \ar[r] & F^nB \ar[r] \ar[d] & B \ar[r] \ar[d] & B' \ar[r] \ar[d] & 0 \\
0 \ar[r] & F^nC \ar[r] & C \ar[r] & C' \ar[r] & 0
}
$$
whose rows are short exact sequences of objects of $\mathcal{A}$.
The proof of (2) is dual. Of course (5) follows from (2).

\medskip\noindent
To prove (4) denote $f : A \to B$ and $g : B \to C$ the given morphisms.
We know that $f(F^p(A)) = \Ker(F^p(B) \to F^p(C))$ by (2) and
$f(A) = \Ker(g)$ by (5). Hence
$f(F^p(A)) =  \Ker(F^p(B) \to F^p(C)) =
\Ker(g) \cap F^p(B) = f(A) \cap F^p(B)$ which proves that
$f$ is strict. The proof that $g$ is strict is dual to this.
\end{proof}











